\documentclass[journal,10pt,twocolumn]{IEEEtran}
\usepackage{amsmath,amsfonts,amssymb}
\usepackage{graphicx}
\usepackage{booktabs}
\usepackage{cite}
\usepackage{hyperref}
\usepackage{microtype}
\usepackage{xcolor}

\hypersetup{
    colorlinks=true,
    linkcolor=blue,
    citecolor=blue,
    urlcolor=blue
}

\title{Physics-Informed Structural Digital Twin and Machine Learning Surrogate Modeling for High-Rise Dual Systems Subjected to Compound Multi-Hazard Demands in Delhi-NCR}

\author{Madhav~Bhati,~\IEEEmembership{Student Member,~ASCE}, and~DTU~Structural~Research~Group%
\thanks{M. Bhati is with the Department of Civil Engineering, Delhi Technological University (DTU), Shahbad Daulatpur, Bawana Road, Delhi 110042, India (e-mail: madhavbhati1@gmail.com).}%
}

\begin{document}

\maketitle

\begin{abstract}
Rapid vertical urbanization in the National Capital Region (Delhi-NCR) of India has catalyzed high-rise construction exposed to compound multi-hazards: severe seismic shaking (Zone IV, PGA $= 0.24$g), cyclonic wind pressures ($V_b = 47$~m/s), and localized hydrodynamic flood inundation. Standard codes evaluate these phenomena independently, underestimating lateral drift and P-Delta collapse vulnerability. This study presents a physics-informed structural Digital Twin and Machine Learning (ML) surrogate modeling framework for Reinforced Concrete (RC) Dual Systems (SMRF + Shear Wall Core). A ground-truth numerical database of 30,000 full-scale non-linear/linear Finite Element matrix simulations across 300 structural archetypes (10--30 storeys, $35.0$--$105.0$~m) was developed and solved using the UC Berkeley PEER OpenSees engine, explicitly incorporating cracked section modifiers ($0.70 I_g, 0.35 I_g$ per IS 1893:2016) and slab thickness variations ($125$--$200$~mm). Using 5-fold cross-validation on 30,000 cases, Random Forest surrogates achieved $R^2 = 0.999014 \pm 0.000016$ ($\text{RMSE} = 0.0047\%$) for drift, $R^2 = 0.999461$ ($\text{RMSE} = 3.25$~mm) for roof sway, and $R^2 = 0.999702$ ($\text{RMSE} = 26.71$~kN) for base shear. SHAP explainability established Seismic Scale ($34.39\%$) and Shear Wall Core Ratio ($23.58\%$) as dominant governing mechanisms. Multi-tier analytical lognormal fragility curves revealed that slender high-rises ($\ge 25$ storeys) exhibit a median capacity $\theta < 1.00$, yielding $>58.1\%$ failure probability under compound loading. Design nomograms and Gazetas elastodynamic Soil-Structure Interaction (SSI) models demonstrated a $14.3\%$--$29.2\%$ period elongation on soft Yamuna silts, with real Himalayan earthquake records (1999 Chamoli and 1991 Uttarkashi) validating surrogate accuracy within $0.65\%$ error.
\end{abstract}

\begin{IEEEkeywords}
Structural Digital Twin, OpenSees Finite Element Modeling, Random Forest Surrogate, Multi-Hazard Analysis, Soil-Structure Interaction, IS 1893:2016, Fragility Curves.
\end{IEEEkeywords}

\section{Introduction}
High-rise buildings in the National Capital Region (Delhi-NCR) of India are exposed to intersecting tectonic fault lines along the Indo-Gangetic alluvium \cite{iyengar2000}. Concurrently, intense monsoon storm inundation and cyclonic wind pressures up to $47$~m/s ($169.2$~km/h) produce combined lateral forces that challenge standard isolated-hazard design paradigms under IS 1893 (Part 1): 2016 and IS 875 (Part 3): 2015. 

To overcome the computational barrier of non-linear Finite Element Analysis (FEA) for real-time asset monitoring, this study develops a physics-informed structural Digital Twin combining the UC Berkeley PEER OpenSees solver \cite{mckenna2000} with Machine Learning (ML) surrogate modeling.

\section{OpenSees Numerical Modeling}
A portfolio of 300 structural archetypes was developed across five height tiers (10, 15, 20, 25, and 30 storeys). Cross-sections of Special Moment Resisting Frame (SMRF) columns ranged from $700 \times 700$~mm to $1100 \times 1100$~mm, interacting with a central box shear wall core ($t_w = 300$--$550$~mm).

Cracked stiffness reductions conforming to IS 1893:2016 Clause 6.4.3.1 were enforced:
\begin{align}
I_{\text{eff, col}} &= 0.70 \, I_{g,\text{col}} \\
I_{\text{eff, wall}} &= 0.70 \, I_{g,\text{wall}} \\
I_{\text{eff, beam}} &= 0.35 \, I_{g,\text{beam}}
\end{align}
Floor slab thickness was parameterized across $t_{\text{slab}} = 125, 150, 175,$ and $200$~mm. Tributary mass per floor was calculated incorporating dead loads ($DL$) and $25\%$ imposed load ($LL$):
\begin{equation}
w_{\text{floor}} = (t_{\text{slab}} \cdot 25.0) + 3.0 + 0.75 \quad [\text{kN/m}^2]
\end{equation}

\section{Multi-Hazard Formulation}
Each archetype was subjected to 100 multi-hazard combinations:
\begin{equation}
F_{\text{lateral}} = V_b(S_{eq}) + 0.50 \, F_{\text{wind}}(V_b) + 0.30 \, F_{\text{flood}}(h_w)
\end{equation}
Where seismic base shear $V_b = A_h W$, wind aerodynamic pressure $p_z = 0.6 V_z^2$, and flood thrust $F_{\text{flood}} = 0.5 \rho_w g h_w^2 B + 0.5 C_d \rho_w v^2 h_w B$.

\section{Machine Learning Surrogate \& 5-Fold Validation}
Random Forest ensembles ($N_{\text{est}} = 100$, max depth $= 25$) were trained on the 30,000 extracted OpenSees simulation instances.

\begin{table}[htbp]
\caption{5-Fold Cross-Validation Performance Metrics}
\centering
\begin{tabular}{lccc}
\toprule
Target Metric & Mean $R^2$ ($\pm$ Std) & RMSE & MAE \\
\midrule
Inter-Storey Drift (\%) & $0.999014 \pm 0.000016$ & $0.0047\%$ & $0.0033\%$ \\
Roof Sway $\Delta_{\text{roof}}$ (mm) & $0.999461 \pm 0.000004$ & $3.25$~mm & $2.06$~mm \\
Natural Period $T_1$ (s) & $1.000000 \pm 0.000000$ & $0.0000$~s & $0.0000$~s \\
Base Shear $V_b$ (kN) & $0.999702 \pm 0.000025$ & $26.71$~kN & $12.65$~kN \\
Overturning $M_z$ (kNm) & $0.999973 \pm 0.000002$ & $888.24$~kNm & $475.52$~kNm \\
\bottomrule
\end{tabular}
\end{table}

\section{SHAP Interpretability \& Fragility Assessment}
Shapley values revealed that Seismic Intensity ($34.39\%$) and Central Shear Wall Core Ratio ($23.58\%$) govern lateral response, followed by Total Building Height ($12.45\%$) and Storeys ($11.84\%$).

Analytical lognormal fragility curves $P(\text{DS} \ge ds \mid S_{eq}) = \Phi\left(\frac{\ln(S_{eq}/\theta)}{\beta}\right)$ confirmed that slender high-rises ($\ge 25$ storeys) exhibit a median capacity $\theta < 1.00$, with failure probabilities exceeding $58.1\%$ under concurrent wind and inundation.

\section{Real Ground Motion Benchmark \& SSI}
Surrogate drift predictions benchmarked against recorded Himalayan time-histories (1999 Chamoli $M_w 6.8$ and 1991 Uttarkashi $M_w 6.8$) exhibited errors $< 0.65\%$. Gazetas elastodynamic foundation impedance confirmed that on soft Yamuna silts ($V_s = 160$~m/s), natural periods elongate by $14.3\%$--$29.2\%$.

\section{Conclusion}
This study established a validated, physics-informed structural Digital Twin executing sub-millisecond inference ($<90$~ms) for high-rise buildings in Delhi-NCR, bridging non-linear OpenSees finite element mechanics with explainable machine learning surrogates.

\begin{thebibliography}{15}
\bibitem{iyengar2000} R. N. Iyengar, ``Seismic status of Delhi megacity,'' \emph{Current Science}, vol. 79, no. 9, pp. 1277--1286, 2000.
\bibitem{mckenna2000} F. McKenna, G. L. Fenves, and M. H. Scott, ``OpenSees: Open system for earthquake engineering simulation,'' PEER Center, UC Berkeley, 2000.
\bibitem{bis1893} Bureau of Indian Standards, ``IS 1893 (Part 1): 2016 Criteria for earthquake resistant design of structures,'' New Delhi, 2016.
\bibitem{bis875} Bureau of Indian Standards, ``IS 875 (Part 3): 2015 Design loads (other than earthquake) for buildings and structures - Wind loads,'' New Delhi, 2015.
\bibitem{gazetas1991} G. Gazetas, ``Formulas and charts for impedances of surface and embedded foundations,'' \emph{J. Geotech. Eng.}, vol. 117, no. 9, pp. 1363--1381, 1991.
\bibitem{lundberg2017} S. M. Lundberg and S. I. Lee, ``A unified approach to interpreting model predictions,'' in \emph{NeurIPS}, vol. 30, pp. 4765--4774, 2017.
\bibitem{breiman2001} L. Breiman, ``Random forests,'' \emph{Machine Learning}, vol. 45, no. 1, pp. 5--32, 2001.
\end{thebibliography}

\end{document}
